
%\clearpage

\section{Conclusion}

In this technical report, we introduce Qwen3, the latest version of the Qwen series. Qwen3 features both thinking mode and non-thinking mode, allowing users to dynamically manage the number of tokens used for complex thinking tasks. The model was pre-trained on an extensive dataset containing 36 trillion tokens, enabling it to understand and generate text in 119 languages and dialects. Through a series of comprehensive evaluations, Qwen3 has shown strong performance across a range of standard benchmarks for both pre-trained and post-trained models, including tasks related to code generation, mathematics, reasoning, and agents.

In the near future, our research will focus on several key areas. We will continue to scale up pretraining by using data that is both higher in quality and more diverse in content. At the same time, we will work on improving model architecture and training methods for the purposes of effective compression, scaling to extremely long contexts, etc.
In addition, we plan to increase computational resources for reinforcement learning, with a particular emphasis on agent-based RL systems that learn from environmental feedback. This will allow us to build agents capable of tackling complex tasks that require inference time scaling.
